\documentclass[11pt]{article}
\usepackage[margin=1in]{geometry}
\usepackage{amsmath,amssymb,amsthm}
\usepackage{graphicx}
\usepackage{microtype}
\usepackage[colorlinks=true,linkcolor=blue,citecolor=blue,urlcolor=blue]{hyperref}

\newtheorem{theorem}{Theorem}
\newtheorem{lemma}{Lemma}
\newtheorem{corollary}{Corollary}
\newcommand{\R}{\mathbb R}
\newcommand{\ND}{\mathcal{ND}[0,1]}
\newcommand{\cC}{\mathfrak c}

\title{Every finite-dimensional nowhere-differentiable subspace\\
extends to a closed continuum-dimensional one}
\author{Candidate full solution for arXiv:2306.01561}
\date{11 August 2026}

\begin{document}
\maketitle

\begin{center}
\fbox{\parbox{0.91\linewidth}{\textbf{Status.} Candidate full proof, likely
valid, requiring human review.  The result affirmatively answers the stated
$(1,\cC)$-spaceability question and proves the stronger
$(n,\cC)$-spaceability statement for every finite positive integer $n$.}}
\end{center}

\section{Source question}

Let $C[0,1]$ be the real Banach space with the supremum norm, and let
\[
 \ND=\{f\in C[0,1]:f\text{ is nowhere differentiable}\}.
\]
Ara\'ujo, Barbosa, Raposo Jr., and Ribeiro state on PDF page 4 of
\cite{ABRR2025} that $\ND$ is known to be spaceable, but ask whether it is
specifically $(1,\cC)$-spaceable.  In other words: does every prescribed
one-dimensional subspace contained in $\ND\cup\{0\}$ lie in a closed
$\cC$-dimensional subspace contained in $\ND\cup\{0\}$?

\begin{figure}[ht]
\centering
\includegraphics[width=\linewidth]{figures/open_problem_crop.png}
\caption{The exact open-question passage in \cite[p.~4]{ABRR2025}, rendered
and cropped from the local source-paper PDF.}
\end{figure}

We use one quantitative theorem of Hencl.  In
\cite[Proposition~2]{Hencl2000}, Hencl proves that for every continuous
increasing function
\(
 \varphi:[0,\infty)\to[0,\infty)
\)
with $\varphi(0)=0$, there are a Cantor set $K\subset[0,1]$ and a linear
isometry
\[
                  T:C(K)\longrightarrow C[0,1]
\]
such that the following holds.  For every $0\ne u\in C(K)$, $z\in(0,1)$,
and $m\in\mathbb N$, the set
\begin{equation}\label{eq:hencl-set}
 D_z(u,m)=\left\{y\in[0,1]\setminus\{z\}:
 \frac{|Tu(y)-Tu(z)|}{\varphi(|y-z|)}>m\right\}
\end{equation}
has symmetric upper density one at $z$.  At $0$ and $1$ it has,
respectively, right and left upper density one.  Hencl also obtains $Tu=u$
on $K$, though that feature is not needed here.

\section{Proof intuition}

A prescribed finite-dimensional subspace $E\subset C[0,1]$ has a common
modulus of continuity: compactness of its unit ball makes the oscillations of
all its unit vectors uniformly small at small scales.  Choose Hencl's
embedding to oscillate much more rapidly than that common modulus.  Then no
fixed element of $E$ can cancel the small-scale oscillations of any nonzero
Hencl function.

Consequently the Hencl range is disjoint from $E$ except at zero, and every
nonzero vector in their algebraic sum is still nowhere differentiable.  The
Hencl range is closed because the embedding is isometric, and adding the
finite-dimensional space $E$ preserves closedness.  This produces the
required closed extension without altering the prescribed directions.

\section{A common modulus for a finite-dimensional subspace}

Let $E$ be a nonzero finite-dimensional subspace of $C[0,1]$.  For
$0\le t\le1$, define
\begin{equation}\label{eq:Omega}
 \Omega_E(t)=\sup\bigl\{|e(x)-e(y)|:
       e\in E,\ \|e\|_\infty\le1,\ x,y\in[0,1],\ |x-y|\le t\bigr\}.
\end{equation}

\begin{lemma}\label{lem:modulus}
The function $\Omega_E$ is increasing, subadditive, and continuous, with
$\Omega_E(0)=0$ and $\Omega_E(t)\to0$ as $t\downarrow0$.  If $E$ contains a
nonconstant function, then for some $c>0$ and all sufficiently small $t>0$,
\[
                         \Omega_E(t)\ge ct.
\]
\end{lemma}

\begin{proof}
The closed unit ball of $E$ is compact in the supremum norm.  A finite
$\varepsilon$-net, together with uniform continuity of its finitely many
members, shows that this ball is uniformly equicontinuous.  Hence
$\Omega_E(t)\to0$ as $t\downarrow0$.

Monotonicity is immediate.  To prove subadditivity, suppose
$|x-y|\le s+t$.  Choose a point between $x$ and $y$ so that the two resulting
distances are at most $s$ and $t$.  Applying the triangle inequality and then
taking the supremum in \eqref{eq:Omega} gives
\(
 \Omega_E(s+t)\le\Omega_E(s)+\Omega_E(t)
\)
when $s+t\le1$.  Thus
\(
 |\Omega_E(s)-\Omega_E(t)|\le\Omega_E(|s-t|),
\)
which, with the limit at zero, proves continuity.

Finally choose $e_0\in E$ with $\|e_0\|_\infty=1$ and points $a<b$ for which
$\delta=|e_0(b)-e_0(a)|>0$.  Partition $[a,b]$ into
$N=\lceil(b-a)/t\rceil$ intervals of length at most $t$.  The triangle
inequality yields
\[
 \delta\le N\Omega_E(t)\le\left(\frac{b-a}{t}+1\right)\Omega_E(t),
\]
and therefore
\(
 \Omega_E(t)\ge \delta t/(b-a+t).
\)
This gives the claimed linear lower bound near zero.
\end{proof}

\section{The finite-extension theorem}

\begin{theorem}\label{thm:main}
For every $n\in\mathbb N$, the set $\ND$ is $(n,\cC)$-spaceable in
$C[0,1]$.  Explicitly, if $E\subset C[0,1]$ is $n$-dimensional and
\[
                         E\setminus\{0\}\subset\ND,
\]
then there is a closed subspace $M\subset C[0,1]$ such that
\[
 E\subset M,\qquad \dim M=\cC,\qquad
 M\setminus\{0\}\subset\ND.
\]
The same proof rules out finite one-sided derivatives at the endpoints.
\end{theorem}

\begin{proof}
Fix such an $E$.  Since every nonzero member of $E$ is nowhere
differentiable, $E$ contains a nonconstant function.  Extend $\Omega_E$
constantly from $[0,1]$ to $[0,\infty)$ and set
\[
                         \varphi(t)=\sqrt{\Omega_E(t)}.
\]
By Lemma~\ref{lem:modulus}, $\varphi$ is continuous and increasing, with
$\varphi(0)=0$.  Moreover,
\begin{equation}\label{eq:two-ratios}
 \frac{\Omega_E(t)}{\varphi(t)}=\sqrt{\Omega_E(t)}\longrightarrow0,
 \qquad
 \frac{\varphi(t)}{t}\longrightarrow\infty
 \quad(t\downarrow0).
\end{equation}
The second assertion follows from $\Omega_E(t)\ge ct$ near zero.

Apply Hencl's proposition to this $\varphi$, and denote the resulting
isometry by $T:C(K)\to C[0,1]$.  Put $Y=T(C(K))$.

We first prove the key perturbation statement.  Fix $0\ne u\in C(K)$,
$e\in E$, and $z\in[0,1]$.  Because $D_z(u,1)$ in
\eqref{eq:hencl-set} has upper density one at $z$ (with the stated one-sided
interpretation at an endpoint), it contains a sequence $y_j\ne z$ converging
to $z$ from the allowed side.  By the definition of $\Omega_E$,
\begin{equation}\label{eq:E-bound}
 |e(y_j)-e(z)|\le\|e\|_\infty\Omega_E(|y_j-z|).
\end{equation}
The first limit in \eqref{eq:two-ratios} shows that, for all sufficiently
large $j$, the right side of \eqref{eq:E-bound} is at most
$\varphi(|y_j-z|)/2$.  On the other hand, membership in $D_z(u,1)$ gives
\(
 |Tu(y_j)-Tu(z)|>\varphi(|y_j-z|).
\)
Hence
\begin{align*}
 |(Tu+e)(y_j)-(Tu+e)(z)|
 &\ge |Tu(y_j)-Tu(z)|-|e(y_j)-e(z)|\\
 &>\frac12\varphi(|y_j-z|).
\end{align*}
After division by $|y_j-z|$, the second limit in
\eqref{eq:two-ratios} shows that these difference quotients are unbounded.
Thus $Tu+e$ has no finite derivative at $z$.  Since $z$ was arbitrary,
\begin{equation}\label{eq:rough-sum}
       0\ne u\in C(K),\ e\in E
       \quad\Longrightarrow\quad Tu+e\in\ND.
\end{equation}

In particular $Y\cap E=\{0\}$: if $Tu=e\in E$ with $u\ne0$, then
\eqref{eq:rough-sum} applied to $Tu-e=0$ would assert that the zero function
is nowhere differentiable, a contradiction.  Now define
\[
                              M=Y+E=Y\oplus E.
\]
The range $Y$ is closed because $T$ is an isometry.  Also $Y+E$ is closed:
if $q:C[0,1]\to C[0,1]/Y$ is the quotient map, then
$q(E)$ is finite-dimensional and hence closed, while
$Y+E=q^{-1}(q(E))$.

Every nonzero element of $M$ has the unique form $Tu+e$.  If $u\ne0$, it
belongs to $\ND$ by \eqref{eq:rough-sum}; if $u=0$, then $e\ne0$ and the
hypothesis on $E$ applies.  Finally $Y$ is isometric to $C(K)$, whose Hamel
dimension is $\cC$, so $\dim M=\cC$.  This proves all assertions.
\end{proof}

\begin{corollary}
The $(1,\cC)$-spaceability question in \cite{ABRR2025} has an affirmative
answer.
\end{corollary}

\section{Verification, scope, and novelty}

The proof is noncomputational.  Its load-bearing points are: compactness of
the unit ball of a finite-dimensional subspace; the two scale comparisons in
\eqref{eq:two-ratios}; the upper-density quantifier in Hencl's Proposition~2;
and closedness of a finite-dimensional extension of a closed subspace.  The
accompanying \texttt{verification.md} records adversarial checks of each.

The theorem concerns real-valued continuous functions, exactly as in the
source passage.  It proves the positive finite-dimensional part.  Ribeiro's
2024 thesis independently proves that $\ND$ is not
$(\alpha,\beta)$-spaceable when $\alpha$ is infinite
\cite[Corollary~2.2]{Ribeiro2024}; together, the two statements give the
finite-versus-infinite classification for prescribed starting dimension.

\paragraph{Bounded novelty check.}
On 11 August 2026, the local run indexes were searched by arXiv id, title,
and the terms ``nowhere differentiable'', ``$(1,\cC)$-spaceable'',
``pointwise spaceability'', ``Hencl'', and ``finite-dimensional extension''.
External exact and close searches used the same phrases and checked the
source/publisher record, Hencl's theorem, Ribeiro's 2024 thesis, and work
indexed through August 2026.  The thesis's pointwise-spaceability result is
for a different set of unbounded integrable functions; for $\ND$ it supplies
only the negative infinite-dimensional obstruction.  The July 2026 preprint
of Lima, Santos, and Ribeiro proves that every prescribed finite-dimensional
subspace extends to a countably infinite-dimensional \emph{algebraic}
subspace \cite{LSR2026}.  It does not claim a closed extension or continuum
dimension, and it does not invoke Hencl's quantitative embedding.  No prior
statement of the theorem above or of this perturbation argument was found.
Novelty confidence is moderate rather than definitive; specialist review of
non-indexed theses and citing literature is recommended.

\begin{thebibliography}{9}
\bibitem{ABRR2025}
G. Ara\'ujo, A. Barbosa, A. Baganha Raposo Jr., and G. Ribeiro,
\emph{Complements of unions: insights on spaceability and applications},
Mathematika \textbf{71} (2025), e70006;
\href{https://arxiv.org/abs/2306.01561}{arXiv:2306.01561}.

\bibitem{Hencl2000}
S. Hencl,
\emph{Isometrical embeddings of separable Banach spaces into the set of
nowhere approximatively differentiable and nowhere H\"older functions},
Proc. Amer. Math. Soc. \textbf{128} (2000), 3505--3511,
\href{https://doi.org/10.1090/S0002-9939-00-05595-7}
{doi:10.1090/S0002-9939-00-05595-7}.

\bibitem{Ribeiro2024}
G. dos Santos Ribeiro,
\emph{A search for linearity in the universe of topological vector spaces},
Ph.D. thesis, Universidade Federal da Para\'iba, 2024,
\href{https://repositorio.ufpb.br/jspui/handle/123456789/38081}
{institutional repository record}.

\bibitem{LSR2026}
F. Lima, M. Santos, and G. Ribeiro,
\emph{On the $(n,\aleph_0)$-lineability of nowhere differentiable functions},
\href{https://arxiv.org/abs/2607.15749}{arXiv:2607.15749} (2026).
\end{thebibliography}

\end{document}
